\documentclass[11pt,reqno]{amsart}

\usepackage{amsmath,amssymb,amsfonts,amsthm,mathtools}
\usepackage{geometry}
\usepackage{hyperref}
\usepackage{booktabs}
\usepackage{array}
\usepackage{enumitem}

\geometry{
  letterpaper,
  margin=1.0in
}

\hypersetup{
  colorlinks=true,
  linkcolor=blue!60!black,
  citecolor=blue!60!black,
  urlcolor=blue!60!black
}

% Theorem Environments
\newtheorem{theorem}{Theorem}[section]
\newtheorem{lemma}[theorem]{Lemma}
\newtheorem{proposition}[theorem]{Proposition}
\newtheorem{corollary}[theorem]{Corollary}

\theoremstyle{definition}
\newtheorem{definition}[theorem]{Definition}

\theoremstyle{remark}
\newtheorem{remark}[theorem]{Remark}

% Custom Mathematical Macros
\newcommand{\CapNull}{\operatorname{Cap}_{\mathrm{null}}}
\newcommand{\CapGlobal}{\operatorname{Cap}_{\mathrm{global}}}
\newcommand{\ResLoc}{\operatorname{Res}_{\mathrm{loc}}}
\newcommand{\FluxBdy}{\mathcal{F}_{\mathrm{boundary}}}
\newcommand{\tr}{\operatorname{tr}}
\newcommand{\PiSz}{\Pi_{\mathrm{Sz}}}
\newcommand{\Id}{\mathbb{I}}

\title[Boundary Filtration, Berezin--Toeplitz POVM Quantization, and Symplectic Capacity]{%
  Paper II: Boundary Filtration of Isotropic Null Rays,\\
  Berezin--Toeplitz POVM Quantization, and Symplectic Capacity%
}

\author{Anthony Curtis Rivera}
\date{\today}

\begin{document}

\begin{abstract}
Building directly upon the machine-verified uniqueness of the critical Cartan Type~IV domain $\mathcal{D}_5 \cong SO_0(5,2)/[SO(5)\times SO(2)]$ and its conformal Shilov horizon $Q_5 \cong (S^4 \times S^1)/\mathbb{Z}_2$ established in Paper~I, we extend the parameter-free deductive chain into the boundary harmonic analysis and Berezin--Toeplitz operator quantization of D5. Using Wallach's discrete analytic continuation of Bergman spaces, Faraut--Kor\'anyi symmetric cone analysis, and Hua's harmonic analysis on classical domains, we isolate the Cartan domain genus p=5 (\texttt{CartanDomainGenus}), the Poisson--Szeg\H{o} kernel homogeneity exponent pPS=5 (\texttt{PoissonSzegoKernelExponent}), and the Cauchy--Szeg\H{o} Hardy reproducing kernel exponent kH=5/2 (\texttt{CauchySzegoHardyExponent}) through formal uniqueness gateways. We then prove that restricting L2(Q5) to the holomorphic Hardy space H2(Q5) via the orthogonal Cauchy--Szeg\H{o} projector $\PiSz$ induces a positive-operator-valued measure (POVM) $\mathcal{P}_E = \PiSz M_{\chi_E} \PiSz$ with strictly positive anti-holomorphic Fourier leakage w(p)=p(1-p)/2>0 (\texttt{AntiHolomorphicFourierLeakage}) for every non-trivial Borel subset EQ5 (0<(E)<1). This machine-certified result refutes sharp projection-valued idempotence (PE2PE) as a theorem of boundary causal holomorphy rather than a physical postulate. Finally, we isolate the complementary trace amplitude areg=1-asing, the quadratic Radon--Nikodym boundary probability weight w=(1-asing)2, and the identical vanishing of the pseudo-Riemannian interval s2(u,v)0 across totally isotropic boundary wavefronts. Every definition and theorem is verified in \textsc{Lean~4} (\texttt{Paper2.lean}) with zero free parameters and zero hardcoded numeric literals.
\end{abstract}

\maketitle
\tableofcontents

%==============================================================================
\section{Introduction and Inheritance from Paper~I}
%==============================================================================

In Paper~I (\texttt{Paper1.lean}), beginning solely from the vanishing generator capacity $\CapNull = 0$ (\texttt{NullCapacity}) along primitive null rays (*Q0), we established via the capability-intersection gateway (\texttt{IntersectionIsolation}) and algebraic conservation gateway (\texttt{UniqueLawSolution}) that:
\begin{enumerate}[label=(\roman*)]
  \item Pointed Lorentzian cone convexity uniquely isolates the timelike signature rank p*=1 (\texttt{TimelikeSignatureRank});
  \item Across the entire Jordan--von Neumann--Wigner and Koecher--Vinberg classification, requiring non-collinear boundary rays (r>1) and a quadratic null boundary variety (r22r) eliminates all higher-rank matrix and exceptional algebras (r3), uniquely isolating the rank-2 Euclidean Jordan spin factor Vn (r1=SymmetricSplitRank=2) and its Kantor--Koecher--Tits Type~IV symmetric tube domain $\mathcal{D}_n \cong SO_0(n,2)/[SO(n)\times SO(2)]$;
  \item Primitive Jordan idempotents cc=c have spectral trace $k^* = \tr(c) = 1$ (\texttt{PrimitiveIdempotentTrace}), scalar component c0=1/2, spatial vector norm squared $\|\mathbf{c}\|^2 = 1/4$, and identically vanishing quadratic norm $Q(c) = c_0^2 - \|\mathbf{c}\|^2 = 0$;
  \item Combining Kronheimer's hyperk\"ahler ALE slice bound (dC=2, \texttt{HyperkahlerSliceComplexDim}) with Helgason's restricted root theory (mlong=1, mshort=r2=3, chiral alternating Weyl order T=|A3|=3) uniquely isolates the critical complex dimension n*=r1+r2=5 (\texttt{CriticalComplexDimension}) and the conformal Shilov horizon $Q_5 \cong (S^d \times S^1)/\mathbb{Z}_2$ of spatial sphere dimension d=n*-p*=4 (\texttt{ShilovHorizonSphereDim}).
\end{enumerate}

In this second installment, we import the isolated variables of \texttt{Paper1.lean} directly into \texttt{Paper2.lean} without introducing a single new assumption or numeric literal. We show that the boundary harmonic analysis, Berezin--Toeplitz quantization, and Born/Radon--Nikodym probability measure on Q5 are uniquely determined by the invariants isolated in Paper~I.

%==============================================================================
\section{Step 9: Isolation of the Cartan Genus and Boundary Reproducing Exponents}
%==============================================================================

\subsection{The Bergman--Hardy Continuation and \texttt{CartanDomainGenus}}

Let H(Dn*) denote the weighted Bergman space of holomorphic functions on the critical Type~IV domain D5 with reproducing kernel K(z,w)=(z,w)-, where (z,w)=1-2zw+(zz)(ww) is the generic Jordan norm polynomial of degree r1=2 (Faraut \& Kor\'anyi 1994). By Wallach's classification of the analytic continuation of the holomorphic discrete series (Wallach 1979), the boundary limit in which the reproducing measure concentrates on the 1-codimensional Shilov horizon $Q_5 \subset \partial\mathcal{D}_5$ is governed by the Cartan domain genus p.

\begin{theorem}[Isolation of \texttt{CartanDomainGenus}]
\label{thm:cartan_genus}
On the critical Type~IV domain Dn* with restricted root multiplicities mlong=RestrictedLongRootMultiplicity=1, mshort=r2=ShortRootMultiplicity=3, and real split rank r1=SymmetricSplitRank=2, let pN denote the twice-weight parameter of the boundary reproducing kernel:
\begin{enumerate}[label=(\alph*)]
  \item \textbf{Capability A (\texttt{CapableOfPositiveBoundaryMeasure}):} By the Faraut--Kor\'anyi--Wallach positivity theorem, positivity of the boundary invariant measure requires:
  \begin{equation}
    p \ge m_{\mathrm{long}} \cdot r_2 + r_1.
  \end{equation}
  For p<mlongr2+r1, the sesquilinear inner product acquires negative-norm states above the continuous Wallach threshold.
  \item \textbf{Capability B (\texttt{CapableOfShilovBoundaryConcentration}):} Concentrating the reproducing measure on the Shilov horizon Qn* rather than the interior bulk of Dn* requires:
  \begin{equation}
    p \le n^* = \mathrm{CriticalComplexDimension}.
  \end{equation}
\end{enumerate}
Across all pN, pgenus=mlongr2+r1=5 (\texttt{CartanDomainGenus}) is the strictly unique integer satisfying both capabilities, and coincides identically with the critical complex dimension (pgenus=n*=5).
\end{theorem}
\begin{proof}
Since mlong=r1-p*=1, we have mlongr2+r1=13+2=5=r1+r2=n*. Thus the lower Wallach positivity threshold and the upper Shilov concentration threshold meet at the single point pgenus=5. By \texttt{IntersectionIsolation} (\texttt{isolation\_CartanDomainGenus}), all p<5 fail Capability~A, all p>5 fail Capability~B, and pgenus=5 is strictly unique.
\end{proof}

\subsection{The Poisson--Szeg\H{o} and Cauchy--Szeg\H{o} Kernel Exponents}

\begin{proposition}[Isolation of \texttt{PoissonSzegoKernelExponent} and \texttt{CauchySzegoHardyExponent}]
\label{prop:szego_exponents}
On the Shilov boundary $Q_{n^*} \cong (S^d \times S^1)/\mathbb{Z}_2$ with spatial sphere dimension d=ShilovHorizonSphereDim=4 and timelike phase rank p*=TimelikeSignatureRank=1:
\begin{enumerate}[label=(\roman*)]
  \item \textbf{Poisson--Szeg\H{o} Homogeneity Exponent (\texttt{PoissonSzegoKernelExponent}):} By Hua's harmonic analysis on Qn* (Hua 1963), the conformal Poisson--Szeg\H{o} kernel P(z,u)=cn|(z,z)|pPS/|(z,u)|2kH has homogeneity degree pPSN uniquely isolated by the conservation law pPSp*pPS-p*=d, yielding pPS=d+p*=5.
  \item \textbf{Cauchy--Szeg\H{o} Hardy Exponent (\texttt{CauchySzegoHardyExponent}):} Because the denominator (z,u) is a polynomial of degree r1=SymmetricSplitRank=2 in the boundary variable uQn*, matching the total boundary measure degree pgenus uniquely isolates the Hardy reproducing exponent kHQ as the solution to:
  \begin{equation}
    r_1 \cdot k_H = p_{\mathrm{genus}} \implies k_H = \frac{p_{\mathrm{genus}}}{r_1} = \frac{5}{2}.
  \end{equation}
\end{enumerate}
\end{proposition}

%==============================================================================
\section{Step 10: Berezin--Toeplitz POVM Quantization and Anti-Holomorphic Fourier Leakage}
%==============================================================================

Let  denote the normalized SO(5)SO(2)-invariant Haar probability measure on Q5, normalized to the primitive idempotent trace (Q5)=k*=PrimitiveIdempotentTraceRat=1. The Hilbert space L2(Q5,d) decomposes orthogonally under the Cauchy--Szeg\H{o} projector $\PiSz : L^2(Q_5) \twoheadrightarrow H^2(Q_5)$ into the holomorphic Hardy space H2(Q5) (consisting of non-negative Fourier modes m0 along the S1 phase generator) and its anti-holomorphic orthogonal complement H2(Q5) (m<0):
\begin{equation}
  L^2(Q_5, d\sigma) = H^2(Q_5) \oplus H^2(Q_5)^\perp.
\end{equation}

For any measurable subset EQ5 with indicator function E (E2=E, E=E) and measure p=(E)Q, the Berezin--Toeplitz quantization (Berezin 1975; Boutet de Monvel \& Guillemin 1981) associates to E the positive-operator-valued measure (POVM) element:
\begin{equation}
  \mathcal{P}_E = \PiSz M_{\chi_E} \PiSz : H^2(Q_5) \to H^2(Q_5).
\end{equation}

\begin{theorem}[Isolation of \texttt{AntiHolomorphicFourierLeakage} and Refutation of Sharp PVM Idempotence]
\label{thm:povm_leakage}
Let EQ5 be a Borel subset of normalized measure p=(E)Q, and let $w_\perp(p) = \|(\Id - \PiSz)\chi_E\|_{L^2(Q_5)}^2$ denote the squared norm of the anti-holomorphic Fourier leakage of E1 into H2(Q5):
\begin{enumerate}[label=(\roman*)]
  \item \textbf{Total Variance Identity:} By indicator idempotence (E2=E) and constant-mode orthogonality, the non-zero S1 Fourier fluctuation $\widetilde{\chi}_E = \chi_E - p$ has total squared norm:
  \begin{equation}
    \|\chi_E - p\|_{L^2(Q_5)}^2 = \int_{Q_5} \chi_E^2\,d\sigma - p^2 = p(k^* - p).
  \end{equation}
  \item \textbf{Conjugation Equipartition Law (\texttt{isolation\_AntiHolomorphicFourierLeakage}):} Because E is real-valued (E=E), complex conjugation interchanges the positive (m>0) and negative (m<0) S1 Fourier sectors with equal L2 weight across the r1=2 non-zero chiralities. Thus w(p) is the strictly unique solution in Q to:
  \begin{equation}
    r_1 \cdot w_\perp(p) = p(k^* - p) \implies w_\perp(p) = \frac{p(k^* - p)}{r_1} = \frac{p(1 - p)}{2}.
  \end{equation}
  \item \textbf{Strict POVM Non-Idempotence (\texttt{sharp\_pvm\_idempotence\_is\_invalid}):} For every non-trivial measurable region $\CapNull < p < k^*$ (0<p<1), the anti-holomorphic leakage satisfies:
  \begin{equation}
    w_\perp(p) > \CapNull \qquad \text{and} \qquad \langle 1, \mathcal{P}_E^2 1 \rangle_{H^2} = p - w_\perp(p) < p = \langle 1, \mathcal{P}_E 1 \rangle_{H^2}.
  \end{equation}
  Consequently, PE2PE: no non-trivial localized boundary indicator can act as a sharp projection-valued measure on the causal Hardy space H2(Q5).
\end{enumerate}
\end{theorem}
\begin{proof}
Since $\PiSz 1 = 1$, we have $\mathcal{P}_E 1 = \PiSz \chi_E = p \cdot 1 + \PiSz(\widetilde{\chi}_E)$ and $\langle 1, \mathcal{P}_E^2 1 \rangle = \|\PiSz \chi_E\|_{L^2}^2 = \|\chi_E\|_{L^2}^2 - \|(\Id - \PiSz)\chi_E\|_{L^2}^2 = p - w_\perp(p)$. For $\CapNull < p < k^*$, both p>0 and k*-p>0, so r1w(p)=p(k*-p)>0. Since r1>0, \texttt{nlinarith} certifies $w_\perp(p) > \CapNull$ and p-w(p)<p.
\end{proof}

%==============================================================================
\section{Step 11: Complementary POVM Amplitude, Radon--Nikodym Weight, and Isotropic Wavefronts}
%==============================================================================

\subsection{Trace Conservation and Quadratic Radon--Nikodym Probability Weight}

When the singular lens link L(3,1)Q5 absorbs a spectral trace amplitude asingQ from the normalized identity trace k*=PrimitiveIdempotentTraceRat=1, the remaining regular horizon amplitude and its induced boundary probability density are uniquely fixed by conservation of trace and rank-r1 quadratic pairing:

\begin{theorem}[Isolation of \texttt{ComplementaryRegularTraceAmplitude} and \texttt{RadonNikodymBoundaryProbabilityWeight}]
\label{thm:radon_nikodym_weight}
For any absorbed singular link trace amplitude asingQ:
\begin{enumerate}[label=(\roman*)]
  \item \textbf{Complementary Regular Trace Amplitude (\texttt{ComplementaryRegularTraceAmplitude}):} By the POVM partition of unity $\mathcal{P}_{\mathrm{sing}} + \mathcal{P}_{\mathrm{reg}} = \Id_{H^2}$, the regular horizon amplitude areg is the strictly unique solution to:
  \begin{equation}
    a_{\mathrm{sing}} + a_{\mathrm{reg}} = k^* \implies a_{\mathrm{reg}}(a_{\mathrm{sing}}) = k^* - a_{\mathrm{sing}} = 1 - a_{\mathrm{sing}}.
  \end{equation}
  \item \textbf{Quadratic Radon--Nikodym Boundary Probability Weight (\texttt{RadonNikodymBoundaryProbabilityWeight}):} Evaluating the sesquilinear Berezin--Toeplitz reproducing pairing across the rank-r1 boundary fiber (r1=2) yields the quadratic conservation law:
  \begin{equation}
    w(a_{\mathrm{sing}}) + a_{\mathrm{sing}}\bigl(r_1 k^* - a_{\mathrm{sing}}\bigr) = (k^*)^{r_1},
  \end{equation}
  whose strictly unique solution is:
  \begin{equation}
    w(a_{\mathrm{sing}}) = \bigl(a_{\mathrm{reg}}(a_{\mathrm{sing}})\bigr)^{r_1} = (1 - a_{\mathrm{sing}})^2.
  \end{equation}
\end{enumerate}
\end{theorem}

\subsection{Identical Vanishing of Pseudo-Riemannian Intervals on Isotropic Wavefronts}

\begin{theorem}[Isolation of \texttt{IsotropicWavefrontIntervalSq} (\texttt{isotropic\_wavefront\_interval\_vanishes})]
\label{thm:isotropic_wavefront}
Let $W_{\mathrm{iso}} = \operatorname{span}\{u,v\} \subset \mathcal{N}$ be a totally isotropic boundary wavefront (2-plane) on which every vector is null: $Q(u) = Q(v) = Q(u+v) = \CapNull$. By the quadratic polarization law:
\begin{equation}
  \Delta s^2(u,v) + Q(u+v) = r_1 Q(u) + r_1 Q(v),
\end{equation}
the squared pseudo-Riemannian interval s2(u,v) (\texttt{IsotropicWavefrontIntervalSq}) is uniquely isolated as:
\begin{equation}
  \Delta s^2(u,v) = \CapNull = 0.
\end{equation}
\end{theorem}

%==============================================================================
\section{Machine-Verified Compiler Summary (\texttt{Paper2.lean})}
%==============================================================================

Table~\ref{tab:paper2_ledger} summarizes the invariants and functional laws isolated in \texttt{Paper2.lean}, all of which are imported directly into \texttt{Paper3.lean} to close the automorphic spectral ledger.

\begin{table}[htbp]
\centering
\caption{Compiler-Verified Invariants and Functional Laws of Paper~II (\texttt{Paper2.lean})}
\label{tab:paper2_ledger}
\renewcommand{\arraystretch}{1.25}
\begin{tabular}{@{}llcc@{}}
\toprule
\textbf{Lean 4 Identifier} & \textbf{Geometric / Operator Meaning} & \textbf{Formula} & \textbf{Value on D5} \\
\midrule
\texttt{CartanDomainGenus} & Wallach--Cartan genus of D5 & mlongr2+r1 & 5 \\
\texttt{PoissonSzegoKernelExponent} & Conformal Poisson--Szeg\H{o} homogeneity degree & d+p* & 5 \\
\texttt{CauchySzegoHardyExponent} & Hardy reproducing kernel exponent on Q5 & pgenus/r1 & 5/2 \\
\texttt{AntiHolomorphicFourierLeakage} & H2(Q5) Fourier leakage of indicator E & p(k*-p)/r1 & p(1-p)/2>0 \\
\texttt{ComplementaryRegularTraceAmplitude} & Depleted regular horizon trace amplitude & k*-asing & 1-asing \\
\texttt{RadonNikodymBoundaryProbabilityWeight} & Quadratic Radon--Nikodym boundary weight & (k*-asing)r1 & (1-asing)2 \\
\texttt{IsotropicWavefrontIntervalSq} & Pseudo-Riemannian interval on isotropic 2-plane & r1Q(u)+r1Q(v)-Q(u+v) & 0 \\
\bottomrule
\end{tabular}
\end{table}

\begin{thebibliography}{99}

\bibitem{Berezin1975}
F.~A.~Berezin, \emph{General concept of quantization}, Comm.\ Math.\ Phys.\ \textbf{40} (1975), 153--174.

\bibitem{BoutetGuillemin1981}
L.~Boutet de Monvel and V.~Guillemin, \emph{The Spectral Theory of Toeplitz Operators}, Annals of Mathematics Studies, Vol.~99, Princeton University Press, Princeton, NJ, 1981.

\bibitem{FarautKoranyi1994}
J.~Faraut and A.~Kor\'anyi, \emph{Analysis on Symmetric Cones}, Oxford Mathematical Monographs, Clarendon Press, Oxford, 1994.

\bibitem{Hua1963}
L.~K.~Hua, \emph{Harmonic Analysis of Functions of Several Complex Variables in the Classical Domains}, Translations of Mathematical Monographs, Vol.~6, American Mathematical Society, Providence, RI, 1963.

\bibitem{Wallach1979}
N.~R.~Wallach, \emph{The analytic continuation of the discrete series. I, II}, Trans.\ Amer.\ Math.\ Soc.\ \textbf{251} (1979), 1--17, 19--37.

\end{thebibliography}

\end{document}
